\documentclass[12pt,a4paper]{article}
\usepackage{ibnp-base}
\usepackage{booktabs}
\usepackage{tabularx}

\ibnptitle{Oficina de Canto Infantil de Guapó}
\ibnpsubject{Projeto Piloto 01 • Impacto Sociocultural}
\ibnplocation{Guapó - GO}
\ibnpdate{Dezembro de 2026}

\begin{document}

% ---------------------------------------------------------
% Cabeçalho / Título Oficial do Projeto
% ---------------------------------------------------------
\begin{center}
    {\color{ibnpNavy}\sffamily\LARGE\bfseries OFICINA DE CANTO INFANTIL DE GUAPÓ}\\[0.3cm]
    {\color{ibnpBlue}\sffamily\large Projeto Piloto de Impacto Sociocultural • Crianças de 8 a 12 anos}\\[0.2cm]
    {\small\textbf{Igreja Batista Nacional da Paz de Guapó (IBNP)} \textbullet{} Escola Social de Guapó}
\end{center}

\vspace{0.5cm}

% ---------------------------------------------------------
% 1. Identificação da Proponente
% ---------------------------------------------------------
\section{Identificação da Proponente}

A proponente é uma organização beneficente e comunitária que atua no fortalecimento de vínculos familiares e sociais no município de Guapó-GO:

\begin{table}[h!]
\centering
\small
\renewcommand{\arraystretch}{1.3}
\begin{tabularx}{\textwidth}{@{}>{\bfseries\color{ibnpNavy}}p{4.2cm}X@{}}
\toprule
Razão Social Oficial: & Igreja Batista Nacional da Paz de Guapó \\
Nome Fantasia / Sigla: & IBNP Guapó \\
CNPJ: & 02.930.019/0001-62 \\
Data de Fundação: & 14 de janeiro de 1999 \\
Sede e Endereço Físico: & Rua Presidente Kennedy, Qd. 21, Lt. 13, Centro, Guapó – GO, CEP 75350-000 \\
Telefone / WhatsApp: & (62) 9870-0089 \\
E-mail Institucional: & \href{mailto:contato@ibnpguapo.org.br}{contato@ibnpguapo.org.br} \\
Portal Oficial: & \href{https://ibnpguapo.org.br}{https://ibnpguapo.org.br} \\
Vínculo Notarial: & 2º Serviço Notarial da Comarca de Guapó - GO \\
Filiação Denominacional: & Convenção Batista Nacional (CBN) e ORMIBAN Goiás \\
Iniciativa Mantida: & Escola Social de Guapó \\
\bottomrule
\end{tabularx}
\end{table}

\vspace{0.3cm}

% ---------------------------------------------------------
% 2. Apresentação do Projeto Piloto
% ---------------------------------------------------------
\section{Apresentação do Projeto Piloto}

A \textbf{Oficina de Canto Infantil de Guapó} configura-se como a iniciativa inaugural de impacto sociocultural estruturada pela IBNP para o município de Guapó - GO. Concebida no âmbito da Escola Social de Guapó, a ação tem como horizonte democratizar o acesso à musicalização, à educação vocal e à expressão artística para crianças no período de férias escolares de dezembro.

A música coral é adotada como instrumento de inclusão, disciplina positiva, convívio coletivo e elevação da autoestima infantojuvenil.

\ibnpheaderbox{Quadro de Parâmetros Oficiais do Projeto}{
\begin{itemize}
    \item \textbf{Título do Projeto:} Oficina de Canto Infantil de Guapó (Projeto Piloto 01).
    \item \textbf{Município de Realização:} Guapó - GO.
    \item \textbf{Público-Alvo:} Crianças de 8 a 12 anos.
    \item \textbf{Capacidade / Atendimento:} 30 vagas totalmente gratuitas.
    \item \textbf{Carga Horária Total:} 6 horas de formação orientada.
    \item \textbf{Formato:} 4 encontros presenciais de 1h30 (90 minutos) cada.
    \item \textbf{Período de Realização:} Mês de Dezembro.
    \item \textbf{Local das Atividades:} Auditório da Sede da IBNP (Rua Presidente Kennedy, Centro, Guapó - GO).
\end{itemize}
}

\vspace{0.3cm}

% ---------------------------------------------------------
% 3. Justificativa de Impacto Sociocultural
% ---------------------------------------------------------
\section{Justificativa de Impacto Sociocultural}

O município de Guapó conta com uma demanda significativa por iniciativas públicas e comunitárias voltadas à ocupação sadia e criativa no mês de dezembro, período de encerramento do calendário escolar regular. A ausência de programas formativos acessíveis intensifica o risco de isolamento digital, ociosidade e exposição a vulnerabilidades sociais.

A educação vocal e a prática do canto em conjunto desempenham papel integrador no desenvolvimento neurológico e emocional infantil. Por meio de exercícios de respiração diafragmática, afinação, consciência rítmica e dicção, as crianças participantes aprimoram sua capacidade de concentração, cognição verbal e sensibilidade estética.

Sob a perspectiva comunitária, a oficina fortalece o senso de pertencimento e solidariedade, promovendo a integração entre participantes de diferentes bairros de Guapó e aproximando pais e responsáveis de um ambiente de acolhimento sadio e seguro.

\vspace{0.3cm}

% ---------------------------------------------------------
% 4. Critérios de Seleção e Atendimento das 30 Vagas
% ---------------------------------------------------------
\section{Critérios de Seleção e Atendimento}

Para assegurar transparência, equidade e foco no público que mais necessita, o processo de oferta das \textbf{30 vagas gratuitas} obedecerá às seguintes diretrizes:

\begin{enumerate}
    \item \textbf{Faixa Etária:} Crianças na faixa etária estrita de 8 a 12 anos completos até o início da oficina.
    \item \textbf{Residência no Município:} Comprovação de residência ou vínculo escolar no município de Guapó - GO.
    \item \textbf{Inscrições Públicas e Gratuitas:} Período público de inscrições presenciais na sede da IBNP e mediante formulário simplificado divulgado nos canais comunitários.
    \item \textbf{Prioridade Social:} Em caso de demanda superior ao número de 30 vagas, terão prioridade crianças oriundas de famílias cadastradas em programas de transferência de renda ou acompanhadas pela assistência comunitária.
    \item \textbf{Autorização dos Responsáveis:} Assinatura obrigatória do Termo de Consentimento, Autorização de Uso de Imagem Institucional e Compromisso de Frequência firmado pelo responsável legal.
    \item \textbf{Gratuidade Integral:} Isenção total de taxas, com fornecimento gratuito de material didático (pastas, letras das músicas e crachás), lanches saudáveis em todos os encontros e certificados.
\end{enumerate}

\newpage

% ---------------------------------------------------------
% 5. Estrutura dos 4 Encontros em Dezembro
% ---------------------------------------------------------
\section{Estrutura dos Encontros em Dezembro}

Os 4 encontros presenciais de 1h30 cada (totalizando 6 horas formativas) estão programados para o mês de dezembro, obedecendo à seguinte sequência pedagógica:

\begin{table}[h!]
\centering
\small
\renewcommand{\arraystretch}{1.3}
\begin{tabularx}{\textwidth}{@{}>{\bfseries}c X c@{}}
\toprule
Encontro & Foco Formativo & Carga Horária \\
\midrule
1º Encontro & Acolhimento, Consciência Corporal e Respiração Diafragmática & 1h30 (90 min) \\
2º Encontro & Afinação Vocal Lúdica, Dicção e Repertório Comunitário & 1h30 (90 min) \\
3º Encontro & Dinâmica Coral Coletiva, Divisão de Vozes e Harmonia & 1h30 (90 min) \\
4º Encontro & Ensaio Geral, Apresentação Pública Final e Certificação & 1h30 (90 min) \\
\midrule
\multicolumn{2}{@{}r}{\textbf{Carga Horária Total Formativa:}} & \textbf{6 horas} \\
\bottomrule
\end{tabularx}
\end{table}

\vspace{0.3cm}

% ---------------------------------------------------------
% 6. Apresentação Pública Final e Certificação
% ---------------------------------------------------------
\section{Apresentação Pública Final e Encerramento}

A culminância do projeto piloto ocorrerá na segunda parte do 4º encontro no mês de dezembro, formatada como uma \textbf{Apresentação Pública Comunitária} no auditório principal da IBNP em Guapó:

\begin{itemize}
    \item \textbf{Mostra Coral:} As 30 crianças participantes interpretarão o repertório ensaiado ao longo da oficina, demonstrando à comunidade o aprendizado técnico e expressivo alcançado.
    \item \textbf{Participação Familiar:} Pais, mães, responsáveis, educadores e vizinhos serão convidados de honra, propiciando um momento de valorização do talento infantojuvenil e integração social.
    \item \textbf{Entrega de Certificados:} Cerimônia solene de outorga de certificados oficiais de participação com chancela da IBNP Guapó para todas as crianças concluintes com aproveitamento e assiduidade.
    \item \textbf{Confraternização Comunitária:} Momento de celebração e lanche comunitário compartilhado entre familiares, equipe docente e voluntários.
\end{itemize}

\vspace{0.5cm}

% ---------------------------------------------------------
% 7. Termo de Encerramento e Assinaturas
% ---------------------------------------------------------
\section{Termo de Encerramento e Assinaturas}

O presente documento formaliza a proposta técnica e institucional da Oficina de Canto Infantil de Guapó, afirmando o compromisso da Igreja Batista Nacional da Paz de Guapó na promoção do desenvolvimento humano e sociocultural da comunidade guapoense.

\vspace{1.5cm}

\begin{center}
    Guapó - GO, Dezembro de 2026.
\end{center}

\vspace{1.5cm}

\noindent
\begin{minipage}[t]{0.48\textwidth}
    \centering
    \rule{0.9\textwidth}{0.5pt}\\[0.1cm]
    \textbf{Pastor Presidente}\\
    Igreja Batista Nacional da Paz de Guapó\\
    CNPJ: 02.930.019/0001-62
\end{minipage}
\hfill
\begin{minipage}[t]{0.48\textwidth}
    \centering
    \rule{0.9\textwidth}{0.5pt}\\[0.1cm]
    \textbf{Coordenação Geral do Projeto}\\
    Escola Social de Guapó\\
    IBNP Guapó
\end{minipage}

\end{document}
